\chapter{The Lefschetz (1,1) Theorem}\label{lefschetz-11}
\chapterstatus{complete}

\section{Introduction}\label{lefschetz-11:section-introduction}

The Lefschetz $(1,1)$-theorem is the one case of the Hodge conjecture that is completely understood, and it holds even with integral coefficients: on a
nonsingular projective variety every integral class of type $(1,1)$ in $H^2$ is the class of a divisor (Theorem~\ref{lefschetz-11:theorem-lefschetz-11}). Its proof
uses the exponential sequence, which exists because the Picard group is the cohomology of a sheaf of abelian groups. With hard Lefschetz it gives the Hodge conjecture
in codimension $n - 1$, hence in dimension at most three. This chapter proves the theorem, draws the consequences, and explains precisely which step has no analogue in
higher codimension. References: \cite[Chapter 1, Section 1]{griffiths-harris}, \cite[Section 11.3]{voisin-hodge-1}, \cite{kodaira-spencer-1953}.

\noindent\textbf{Prerequisites.} Chapter~\ref{hodge-conjecture} (The Hodge Conjecture) and Chapter~\ref{gaga} (Analytic Spaces and GAGA).

Throughout, $X$ is a compact K\"ahler manifold of dimension $n$; from Section~\ref{lefschetz-11:section-projective} on it is projective.

\section{The theorem for compact K\"ahler manifolds}\label{lefschetz-11:section-kahler}

\begin{theorem}[Lefschetz $(1,1)$, K\"ahler form]\label{lefschetz-11:theorem-kahler}
Let $X$ be a compact K\"ahler manifold. The image of the first Chern class $c_1 \colon \Pic(X) \to H^2(X; \ZZ)$ is exactly the group of integral Hodge classes
$\mathrm{Hdg}^1(X; \ZZ)$, the classes whose image in $H^2(X; \CC)$ is of type $(1,1)$.
\end{theorem}

\begin{proof}
The Chern class of a holomorphic line bundle is represented by $\frac{i}{2\pi}F$, a closed real $(1,1)$-form (Theorem~\ref{holomorphic-bundles:theorem-chern-connection}), so
$c_1(\Pic(X)) \subseteq \mathrm{Hdg}^1(X; \ZZ)$. Conversely, let $a \in H^2(X; \ZZ)$ have image of type $(1,1)$. By the \emph{key lemma} (Lemma~\ref{kodaira:lemma-integral-11}), the map
$H^2(X; \ZZ) \to H^2(X, \mathcal{O}_X)$ induced by $\ZZ_X \subseteq \mathcal{O}_X$ is the composite of $H^2(X; \ZZ) \to H^2(X; \CC)$ with the projection to $H^{0,2}(X) \cong H^2(X, \mathcal{O}_X)$.
So $a$ maps to $0$ in $H^2(X, \mathcal{O}_X)$. In the long exact sequence of the exponential sequence $0 \to \ZZ \to \mathcal{O}_X \to \mathcal{O}_X^* \to 0$
(Theorem~\ref{holomorphic-bundles:theorem-exponential}),
\[
H^1(X, \mathcal{O}_X^*) = \Pic(X) \xrightarrow{c_1} H^2(X; \ZZ) \to H^2(X, \mathcal{O}_X),
\]
so $a$ is in the image of $c_1$.
\end{proof}

\begin{corollary}\label{lefschetz-11:corollary-neron-severi}
$\mathrm{NS}(X) := c_1(\Pic(X)) = \mathrm{Hdg}^1(X; \ZZ)$, and $\Pic(X)$ is an extension $0 \to \Pic^0(X) \to \Pic(X) \to \mathrm{NS}(X) \to 0$ of a finitely generated abelian group by the complex torus
$\Pic^0(X) = H^1(X, \mathcal{O}_X)/H^1(X; \ZZ)$ of dimension $h^{0,1} = b_1/2$. The \emph{Picard number} $\rho(X) = \operatorname{rank}\mathrm{NS}(X)$ satisfies $\rho(X) \leq h^{1,1}(X)$.
\end{corollary}

\begin{proof}
The first statement is Theorem~\ref{lefschetz-11:theorem-kahler}, and the extension is the exponential sequence. $\mathrm{NS}(X) \subseteq H^2(X; \ZZ)$ is finitely generated, and
$\Pic^0(X)$ is a torus by Remark~\ref{holomorphic-bundles:remark-picard-variety} and Exercise~\ref{hodge-decomposition:exercise-h1-real}. Finally $\mathrm{NS}(X) \otimes \RR$ injects into the real
classes of type $(1,1)$, which form a real vector space of dimension $h^{1,1}$ (a real structure on $H^{1,1}$, since $\overline{H^{1,1}} = H^{1,1}$).
\end{proof}

\section{The projective case: classes of divisors}\label{lefschetz-11:section-projective}

\begin{theorem}[Lefschetz $(1,1)$]\label{lefschetz-11:theorem-lefschetz-11}
Let $X$ be a nonsingular complex projective variety. Every class in $\mathrm{Hdg}^1(X; \ZZ)$ is the class of a divisor: the cycle class map $\CH^1(X) \to \mathrm{Hdg}^1(X; \ZZ)$ is surjective. In particular
$\mathrm{HC}(X, 1)$ holds, even integrally.
\end{theorem}

\begin{proof}
By Theorem~\ref{lefschetz-11:theorem-kahler}, a class $a \in \mathrm{Hdg}^1(X; \ZZ)$ is $c_1(L)$ for a holomorphic line bundle $L$ on $X^{\mathrm{an}}$. By GAGA
(Corollary~\ref{gaga:corollary-consequences}(2)), $L$ is the analytification of an algebraic line bundle, still denoted $L$. Let $H$ be a hyperplane section. By Serre's theorem
(Theorem~\ref{coherent-sheaves:theorem-serre-generation}), $L(m) = L \otimes \mathcal{O}(mH)$ is globally generated for $m \gg 0$, so it has a nonzero section, whose divisor $D_1$ is effective with
$\mathcal{O}(D_1) \cong L(m)$ (Proposition~\ref{divisors:proposition-cartier-picard}). Likewise $\mathcal{O}(mH)$ has the divisor $D_2 = mH$. So $L \cong \mathcal{O}(D_1 - D_2)$, and
$a = c_1(L) = \mathrm{cl}(D_1) - \mathrm{cl}(D_2)$ by Example~\ref{cycle-class:example-classes}(2).
\end{proof}

\begin{corollary}\label{lefschetz-11:corollary-codimension-n-1}
For a nonsingular projective variety $X$ of dimension $n$, $\mathrm{HC}(X, n - 1)$ holds. Consequently the Hodge conjecture holds for all nonsingular projective varieties of
dimension at most $3$.
\end{corollary}

\begin{proof}
$\mathrm{HC}(X, 1)$ holds by Theorem~\ref{lefschetz-11:theorem-lefschetz-11}, and implies $\mathrm{HC}(X, n - 1)$ for $n \geq 2$ by Proposition~\ref{hodge-conjecture:proposition-lefschetz-reduction}.
Together with the trivial cases $p = 0$ and $p = n$ (Proposition~\ref{hodge-conjecture:proposition-trivial-degrees}), this covers all $p$ when $n \leq 3$.
\end{proof}

\begin{remark}\label{lefschetz-11:remark-integral-curves}
For curves on a threefold, the integral version fails in general: Proposition~\ref{hodge-conjecture:proposition-lefschetz-reduction} only gives that some multiple of an integral Hodge
class in $H^4$ is algebraic, since $L^{n-2}$ need not be an isomorphism over $\ZZ$. Koll\'ar's counterexamples to the integral Hodge conjecture are integral classes in $H^4$ of threefolds
(Chapter~\ref{failures}). On the other hand, Voisin proved the integral Hodge conjecture for curve classes on uniruled threefolds and on Calabi--Yau threefolds \cite{voisin-2006}.
\end{remark}

\begin{example}\label{lefschetz-11:example-neron-severi}
\begin{enumerate}
\item For $X \subseteq \PP^3$ a nonsingular quartic surface, $\mathrm{NS}(X)$ is a sublattice of $H^2(X; \ZZ) \cong \ZZ^{22}$ of rank $\rho \in [1, 20]$, and $\rho = 20$ is attained (for instance by the Fermat quartic).
\item For $E \times E'$, $\mathrm{NS} \cong \ZZ^2 \oplus \Hom(E, E')$ (Exercise~\ref{abelian-varieties:exercise-ns-product}), which by Theorem~\ref{lefschetz-11:theorem-lefschetz-11} is also the
lattice of integral classes of type $(1,1)$, a statement about Hodge structures (Example~\ref{hodge-conjecture:example-elliptic-product}).
\end{enumerate}
\end{example}

\begin{theorem}[Noether--Lefschetz]\label{lefschetz-11:theorem-noether-lefschetz}
For $d \geq 4$, a very general nonsingular surface of degree $d$ in $\PP^3$ has $\Pic(X) = \ZZ\cdot\mathcal{O}_X(1)$, so $\rho(X) = 1$.
\end{theorem}

This is proved in \cite{voisin-hodge-2} and \cite{griffiths-harris-noether-lefschetz}. The proof shows, with the infinitesimal variation of Hodge structure (Chapter~\ref{vhs}), that
the locus of surfaces on which a given integral class stays of type $(1,1)$ is a proper subvariety of the parameter space, unless the class is a multiple of the hyperplane class; a
very general surface lies outside all these countably many loci. For $d \leq 3$ the statement is false: quadrics contain two families of lines, and cubic surfaces contain $27$ lines.

\section{Why the proof does not generalise}\label{lefschetz-11:section-why-not}

\begin{remark}\label{lefschetz-11:remark-why-not}
The proof of Theorem~\ref{lefschetz-11:theorem-lefschetz-11} has three steps, and none generalises to codimension $p \geq 2$.
\begin{enumerate}
\item \emph{A cohomological object parametrising cycles.} Divisors modulo linear equivalence are line bundles, $\Pic(X) = H^1(X, \mathcal{O}_X^*)$, the cohomology of a sheaf of abelian
groups, and the exponential sequence relates it to $H^2(X; \ZZ)$. For codimension $p$ the Chow group $\CH^p(X) \cong H^p(X, \mathcal{K}_p)$ (Bloch's formula) involves a sheaf of Quillen
$K$-groups, which has no comparison with singular cohomology that is exact in the needed way.
\item \emph{The key lemma.} The obstruction to lifting an integral class is its image in $H^2(X, \mathcal{O}_X) = H^{0,2}$, a coherent cohomology group, and it vanishes exactly for type $(1,1)$. For
$p \geq 2$, the analogous vanishing of the components in $H^{p-1,p+1}, \ldots, H^{0,2p}$ has no sheaf-theoretic interpretation producing an object like a line bundle. Deligne cohomology
$H^{2p}_{\mathcal{D}}(X; \ZZ(p))$ fits into an exact sequence $0 \to J^p(X) \to H^{2p}_{\mathcal{D}}(X; \ZZ(p)) \to \mathrm{Hdg}^p(X; \ZZ) \to 0$, so every integral Hodge class lifts to Deligne
cohomology (Chapter~\ref{intermediate-jacobians}); the missing step is to lift further to cycles.
\item \emph{Sections of line bundles.} A line bundle twisted by a high multiple of an ample bundle has sections, whose zero loci are divisors. In codimension $\geq 2$ there is no
analogue: a vector bundle of rank $p$ with enough sections gives cycles as zero loci of sections, but the Chern classes of such bundles need not span the Hodge classes, and producing a
bundle with a prescribed Chern class is the vector bundle form of the conjecture itself (Proposition~\ref{hodge-conjecture:proposition-equivalent-forms}(3)).
\end{enumerate}
\end{remark}

\section{Exercises}\label{lefschetz-11:section-exercises}

\begin{exercise}\label{lefschetz-11:exercise-picard-number-torus}
Let $T$ be a complex torus of dimension $2$ with $\rho(T) = 0$. Show that $T$ carries no nonconstant meromorphic functions and no curves.
\end{exercise}

\begin{solution}
A curve $C$ would give a line bundle $\mathcal{O}(C)$ with $c_1(\mathcal{O}(C)) = \mathrm{cl}(C) \neq 0$, since $\int_T\mathrm{cl}(C) \wedge \omega = \int_C\omega > 0$ for a K\"ahler form $\omega$
(Proposition~\ref{kahler:proposition-kahler-class-powers}). So $\mathrm{NS}(T) \neq 0$, a contradiction. A nonconstant meromorphic function $f$ has a nonzero polar divisor, which is a union of
curves.
\end{solution}

\begin{exercise}\label{lefschetz-11:exercise-hodge-classes-surface}
Let $S$ be a nonsingular projective surface with $p_g(S) = h^{2,0}(S) = 0$. Show that every class in $H^2(S; \ZZ)$ is the class of a divisor, and that $\rho(S) = b_2(S)$.
\end{exercise}

\begin{solution}
If $h^{2,0} = h^{0,2} = 0$, then $H^2(S; \CC) = H^{1,1}$, so every integral class is an integral Hodge class, hence algebraic by Theorem~\ref{lefschetz-11:theorem-lefschetz-11}. So
$\mathrm{NS}(S) = H^2(S; \ZZ)$ and $\rho = b_2$.
\end{solution}

\begin{exercise}\label{lefschetz-11:exercise-cubic-threefold}
Show that on a nonsingular cubic threefold $X \subseteq \PP^4$ the Hodge conjecture holds, and compute $\mathrm{Hdg}^1$ and $\mathrm{Hdg}^2$.
\end{exercise}

\begin{solution}
By Corollary~\ref{lefschetz-11:corollary-codimension-n-1} the conjecture holds in dimension $3$. By the Lefschetz hyperplane theorem and Poincar\'e duality, $H^2 = \QQ h$ and $H^4 = \QQ h^2/3$
(the class of a line, since $\int_Xh^3 = 3$), both algebraic, so $\mathrm{Hdg}^1 = H^2(X; \QQ)$ and $\mathrm{Hdg}^2 = H^4(X; \QQ)$. The interesting group $H^3$, of dimension $10$ with $h^{2,1} = 5$, has
no Hodge classes, being of odd degree.
\end{solution}
